%! Skills Focus: Selecting an Appropriate Inference Procedure for Categorical Data — Unit 8, Lesson 8.7 — Homework
\documentclass[10pt]{article}
\usepackage{apstats-article}
\usepackage{apstats-boxes}

\begin{document}

\pageheader{Unit 8, Lesson 8.7}{Homework}
\namedateperiod

\vspace{0.10in}

\begin{enumerate}

\item \textbf{Identify the Procedure.} \; For each scenario, name the correct inference
      procedure. No computation required.

\begin{enumerate}[label=\alph*.]
  \item A meteorologist records weather outcomes (Sunny, Cloudy, Rainy, Snowy) for a random
        sample of 365 days and wants to test whether the four outcomes are equally likely.

        Procedure: \blank{9cm}

  \item A university admissions office randomly samples 80 applicants from in-state and 80
        applicants from out-of-state and records whether each was admitted (yes or no). They
        want to estimate the difference in admission rates.

        Procedure: \blank{9cm}

  \item A market researcher takes one random sample of 500 shoppers and records each shopper's
        age group (Under 30, 30--50, Over 50) and preferred shopping channel (Online, In-Store,
        Both). She wants to know whether shopping channel preference is associated with age group.

        Procedure: \blank{9cm}

  \item A nutrition researcher randomly samples students from four different dining halls and
        records each student's beverage choice (Water, Soda, Juice, Other). She wants to test
        whether the distribution of beverage choice is the same across all four dining halls.

        Procedure: \blank{9cm}
\end{enumerate}

\vspace{0.10in}

\item \textbf{State and Plan.} \;
      A national survey organization claims that 15\% of adults read a physical newspaper daily.
      A researcher suspects the true proportion is lower. She surveys a random sample of 400 adults
      and finds that 48 of them read a physical newspaper daily.

      \begin{enumerate}[label=\alph*.]
        \item \textbf{State:} Procedure: \blank{7cm}\\
              $H_0$: \blank{5cm} \quad $H_a$: \blank{5cm}
        \item \textbf{Plan:} Verify conditions (use numbers):
              \begin{itemize}
                \item Random: \writeline
                \item 10\%: \writeline
                \item Large Counts ($np_0$ and $n(1-p_0)$ each $\ge 10$): \writeline
              \end{itemize}
      \end{enumerate}

\vspace{0.10in}

\item \textbf{AP-Style Multiple Choice.} \;
      \textit{A researcher takes a random sample of 250 households from a large city and records
      each household's neighborhood type (Urban, Suburban, Rural) and internet connection type
      (Broadband, Satellite, None). He wants to test whether there is an association between
      neighborhood type and internet connection type. Which procedure is most appropriate?}

      \begin{enumerate}[label=(\Alph*)]
        \item Chi-square goodness-of-fit test
        \item Chi-square test for homogeneity
        \item Chi-square test for independence
        \item Two-sample $z$-test for $p_1 - p_2$
      \end{enumerate}

\vspace{0.10in}

\item \textbf{Compare Two Procedures.} \;
      A researcher studies whether political party affiliation (Democrat, Republican, Independent)
      differs between male and female voters. She randomly samples 200 male voters and 200 female
      voters separately.

      \begin{enumerate}[label=\alph*.]
        \item Name the most appropriate procedure and explain your reasoning in two sentences.
              \writeline \writeline

        \item A classmate argues that a two-sample $z$-test for $p_1 - p_2$ should be used
              because there are only two groups. Identify the error in this reasoning.
              \writeline \writeline

        \item Write $H_0$ for the correct procedure.
              $H_0$: \writeline
      \end{enumerate}

\vspace{0.10in}

\item \textbf{Conditions and Hypotheses.} \;
      A high school counselor randomly selects 40 freshmen, 40 sophomores, 40 juniors, and
      40 seniors. She records whether each student has participated in at least one extracurricular
      activity this year (yes or no). She wants to test whether the proportion who participate
      differs across grade levels.

      \begin{enumerate}[label=\alph*.]
        \item Name the procedure: \blank{9cm}
        \item $H_0$: \writeline
        \item $H_a$: \writeline
        \item State the key condition that must be checked (beyond Random and 10\%): \writeline
      \end{enumerate}

\end{enumerate}

\begin{extensionbox}
\textbf{Extension: Why Can't Chi-Square Produce a Confidence Interval?} (optional)

A two-sample $z$-test for $p_1 - p_2$ has a corresponding confidence interval: the two-sample
$z$-interval for $p_1 - p_2$. Explain in 2--3 sentences why chi-square procedures (GOF,
homogeneity, independence) do not have a corresponding confidence interval. (\textit{Hint:
think about what parameter a chi-square test is making a claim about.})
\end{extensionbox}

\begin{reflectionbox}
\textbf{Preview: Unit 8 Sample Test}

In the next class, you will complete the Unit 8 Sample Test, which covers all five categorical
inference procedures. Review the decision table from today's guided notes and be ready to select
a procedure, state hypotheses, verify conditions, and communicate a conclusion --- all under
AP-exam conditions.
\end{reflectionbox}

\end{document}
